\chapter{Introducción}
\label{cap:introduccion}

\section{Motivación}
El cáncer de mama es la neoplasia más frecuente en mujeres y una de las
principales causas de mortalidad por cáncer a nivel mundial. La detección precoz
mediante programas de cribado con mamografía reduce de forma significativa la
mortalidad, ya que permite identificar lesiones en estadios tempranos y más
tratables. Sin embargo, la interpretación de mamografías es una tarea compleja,
sujeta a variabilidad inter-observador y a una carga de trabajo elevada para los
radiólogos.

En este contexto, los sistemas de diagnóstico asistido por computador
(\textit{Computer-Aided Diagnosis}, CAD) y, más recientemente, las técnicas de
aprendizaje profundo (\textit{deep learning}) se han convertido en herramientas
de gran interés. Trabajos como el de Yala \textit{et al.}~\cite{yala2019} han
demostrado que los modelos basados en mamografía pueden mejorar la predicción del
riesgo de cáncer de mama. No obstante, estos modelos presentan limitaciones
importantes: son difíciles de interpretar, requieren grandes cantidades de datos
etiquetados y pueden ser sensibles al ruido y a las condiciones de adquisición.

El \textbf{Análisis Topológico de Datos} (TDA) surge como un enfoque
complementario que caracteriza la \emph{forma} intrínseca de los datos mediante
herramientas de la topología algebraica. Su técnica principal, la
\emph{homología persistente}, es robusta frente al ruido y proporciona
descriptores interpretables geométricamente, lo que la hace especialmente
atractiva para el análisis de imagen médica.

\section{Contexto del problema}
Una mamografía digital es una imagen de rayos X de la mama que revela
estructuras como el tejido glandular, los conductos, las calcificaciones y
posibles masas. La densidad mamaria (clasificada habitualmente según el sistema
BI-RADS) y la morfología de las lesiones son factores clave en el diagnóstico.
El conjunto de datos público \mbox{CBIS-DDSM}~\cite{lee2017cbisddsm} proporciona
mamografías digitalizadas con anotaciones de lesiones (masas y
microcalcificaciones), regiones de interés (ROI) segmentadas y etiquetas de
patología (benigno/maligno), lo que lo convierte en un banco de pruebas
adecuado para este trabajo.

\section{Objetivos y alcance (resumen)}
El objetivo general de este TFG es \textbf{estudiar y evaluar la aplicación del
Análisis Topológico de Datos a la clasificación y visualización de mamografías},
comparándolo y combinándolo con técnicas de aprendizaje profundo. Los objetivos
específicos se detallan en el Capítulo~\ref{cap:objetivos}.

\section{Estructura de la memoria}
El resto de la memoria se organiza como sigue:
\begin{itemize}
  \item El \textbf{Capítulo~\ref{cap:objetivos}} formula los objetivos y la
        metodología general del trabajo.
  \item El \textbf{Capítulo~\ref{cap:estado-arte}} revisa el estado del arte en
        análisis de mamografías y en TDA aplicado a imagen médica.
  \item El \textbf{Capítulo~\ref{cap:fundamentos}} introduce los fundamentos
        teóricos de la homología persistente y del aprendizaje profundo.
  \item El \textbf{Capítulo~\ref{cap:metodologia}} describe la metodología
        propuesta: datos, preprocesado, extracción de descriptores y modelos.
  \item El \textbf{Capítulo~\ref{cap:implementacion}} detalla la implementación
        y las herramientas utilizadas.
  \item El \textbf{Capítulo~\ref{cap:resultados}} presenta y discute los
        resultados experimentales.
  \item El \textbf{Capítulo~\ref{cap:conclusiones}} recoge las conclusiones y
        las líneas de trabajo futuro.
\end{itemize}
